\section{Modelagem do Problema}
\label{sec:modelagem}

Para resolver o problema por meio de algoritmos de busca, ele é modelado como
um Problema em Espaço de Estados \parencite{poole2010}. Isso exige a definição
de um conjunto de estados, um estado inicial, um conjunto de ações disponíveis
para cada estado, uma função sucessora que produz os estados alcançáveis a
partir de um estado e uma ação, e uma função objetivo que determina se um
estado é uma solução.

\subsection{Representação dos estados}

Cada estado $s_i \in S$ é representado como uma tripla $s_i = (A, B, \alpha)$,
sendo:
\begin{itemize}
	\item $A$ o conjunto de pessoas do lado de origem da ponte;
	\item $B$ o conjunto de pessoas do lado de destino da ponte;
	\item $\alpha$ uma variável booleana que indica se a tocha está do lado
	      $A$ (verdadeira) ou do lado $B$ (falsa).
\end{itemize}

Na implementação, essa tripla é acrescida de um quarto valor correspondente à
estimativa heurística do estado, utilizado apenas pela busca A* para
priorizar a expansão de estados mais promissores; esse valor não faz parte da
definição formal do estado e não influencia as demais buscas.

\subsection{Estado inicial e estado objetivo}

O estado inicial $s_0 = (P, \{\}, \text{true})$ representa a situação em que
todas as pessoas $P$ e a tocha estão do lado de origem, e o lado de destino
está vazio.

Um estado $s_i$ é considerado objetivo quando o lado de origem está vazio, o
lado de destino não está vazio e a tocha está do lado de destino, isto é,
$s_i = (\{\}, P, \text{false})$. É essa condição que o programa verifica para
determinar se a busca alcançou uma solução.

\subsection{Ações possíveis}

Em cada estado, a tocha se encontra em um dos dois lados da ponte, e apenas
pessoas desse mesmo lado podem se mover. As ações possíveis são a travessia de
uma ou duas pessoas do lado onde está a tocha para o lado oposto -- nunca mais
de duas, dada a capacidade da ponte. Chamamos essas ações de
$q \in \{\mathit{ir}, \mathit{voltar}\}$: $\mathit{ir}$ quando o movimento
ocorre de $A$ para $B$ (tocha em $A$), e $\mathit{voltar}$ quando ocorre de
$B$ para $A$ (tocha em $B$). Como a tocha é necessária para enxergar a ponte,
toda travessia -- de ida ou de volta -- desloca a tocha junto com as pessoas
que se movem, alternando seu lado a cada ação.

\subsection{Função sucessora}

A função sucessora tem assinatura $F(s_i, p_i, q_i)$, em que $q_i$ é o tipo de
ação ($\mathit{ir}$ ou $\mathit{voltar}$) e $p_i \subseteq P$ é o conjunto de
uma ou duas pessoas que atravessam a ponte a partir do estado $s_i$, sempre
pertencentes ao lado de onde a tocha parte. A função retorna o estado
resultante de transferir as pessoas de $p_i$ entre os lados $A$ e $B$ e de
inverter a posição da tocha:

\begin{equation}
	\begin{cases}
		F(s_i, p_i, q_i=\mathit{ir}) = (A \setminus p_i,\; B \cup p_i,\; \text{false})
		\\
		F(s_i, p_i, q_i=\mathit{voltar}) = (A \cup p_i,\; B \setminus p_i,\; \text{true})
	\end{cases}
\end{equation}

Dado um estado, todos os seus sucessores são obtidos aplicando $F$ a cada
combinação possível de uma ou duas pessoas do lado onde está a tocha. Como a
mesma ação pode levar de volta a um estado já visitado (por exemplo, uma
pessoa que atravessa e depois retorna sozinha), a função sucessora não evita
ciclos: essa responsabilidade fica a cargo de cada algoritmo de busca, que
mantém um conjunto de estados já visitados para não reexpandi-los.

\subsection{Custos das transições}

O custo de uma transição corresponde ao tempo gasto na travessia, que é
determinado pela pessoa mais lenta entre as que atravessam juntas -- já que
todas as pessoas do grupo cruzam a ponte no mesmo ritmo, limitado por quem
demora mais. Sendo $t(p)$ o tempo de travessia da pessoa $p$ e $M = p_i$ o
conjunto de pessoas que se movem em uma ação, o custo dessa transição é dado
por:

\[
	c(M) = \max_{p \in M} t(p)
\]

\subsection{Representação do problema como grafo}

O espaço de estados pode ser visto como um grafo direcionado e ponderado, no
qual cada vértice corresponde a um estado $s_i \in S$ e cada aresta
corresponde a uma ação válida $F(s_i, p_i, q_i)$ que leva de um estado a
outro. O peso de cada aresta é o custo $c(M)$ da travessia associada. O grafo
é direcionado porque as ações de ida e de volta não são simétricas -- a tocha
precisa estar no lado de origem da ação -- e pode conter ciclos, já que é
possível retornar a um estado previamente visitado. Resolver o problema
corresponde, portanto, a encontrar um caminho no grafo entre o vértice que
representa o estado inicial $s_0$ e algum vértice que satisfaça a condição de
estado objetivo, sendo o caminho de menor custo total a solução ótima.
